\section{Evidence State Manager (ESM): Policy and Invariants}
\label{sec:c1_esm}

\subsection{ESM objective}
Given a fully PV-scored evidence pool $\Pool$, ESM constructs a disjoint partition:
\begin{equation}
\Pool = \Active\;\dot\cup\;\Susp\;\dot\cup\;\Counter.
\end{equation}

Semantics:
\begin{itemize}
\item $\Active$: evidence consumed by downstream reasoner.
\item $\Susp$: retained low-priority evidence for potential recovery.
\item $\Counter$: contradiction-leaning evidence protected for safety diagnostics.
\end{itemize}

\subsection{Configuration defaults}
\begin{table}[H]
\centering
\begin{tabular}{@{}p{0.24\linewidth}p{0.16\linewidth}p{0.54\linewidth}@{}}
\toprule
Field & Default & Description \\
\midrule
\texttt{min\_A} & 5 & Minimum Active size target (when feasible). \\
\texttt{max\_A} & 20 & Maximum Active size cap. \\
\texttt{max\_C} & 5 & Maximum Counter size cap. \\
\texttt{contra\_tau} & 0.6 & Minimum contradiction probability for Counter candidacy. \\
\bottomrule
\end{tabular}
\caption{\texttt{ESMConfig} defaults in \texttt{src/component1/esm.py}.}
\label{tab:c1_esm_defaults}
\end{table}

\subsection{Deterministic selection policy}
Let each item $e\in\Pool$ have PV fields available.

\paragraph{Step 1: Counter candidates.}
\begin{equation}
\Counter_{\text{cand}} = \{e\in\Pool \mid \Con(e)\ge\tau\}, \quad \tau=\texttt{contra\_tau}.
\end{equation}
Sort $\Counter_{\text{cand}}$ by $\Con(e)$ descending.

\paragraph{Step 2: Active-starvation guard for Counter cap.}
To preserve Active feasibility, code computes:
\begin{align}
\texttt{max\_C\_safe} &= \max(0, |\Pool|-\texttt{min\_A}), \\
\texttt{actual\_max\_C} &= \min(\texttt{max\_C}, \texttt{max\_C\_safe}).
\end{align}
Then:
\begin{equation}
\Counter = \text{top-}\texttt{actual\_max\_C}(\Counter_{\text{cand}}).
\end{equation}

\paragraph{Step 3: Active selection on remaining pool.}
Let $R=\Pool\setminus\Counter$, sorted by $\Rel(e)$ descending.
If $|R|<\texttt{min\_A}$, set $\Active=R$ and $\Susp=\varnothing$.
Else:
\begin{align}
\texttt{num\_active} &= \max(\texttt{min\_A}, \min(\texttt{max\_A}, |R|)), \\
\Active &= R[:\texttt{num\_active}], \\
\Susp &= R[\texttt{num\_active}:].
\end{align}

\subsection{Invariants enforced by assertions}
\begin{enumerate}
\item Active bound: $|\Active|\le\texttt{max\_A}$.
\item Counter bound: $|\Counter|\le\texttt{max\_C}$.
\item Feasibility guard: if $|R|\ge\texttt{min\_A}$ then $|\Active|\ge\texttt{min\_A}$.
\item Completeness: $|\Active|+|\Susp|+|\Counter|=|\Pool|$.
\item Disjointness: no evidence id appears in more than one set.
\end{enumerate}

\subsection{Edge-case behavior}
\paragraph{Empty pool.}
Returns \texttt{\{"A":[], "S":[], "C":[]\}} with warning log.

\paragraph{Pool smaller than \texttt{min\_A}.}
All non-counter remainder goes to Active, and Suspended remains empty.
This is an intentional fallback to avoid artificial starvation.

\paragraph{High contradiction density.}
Even if many items satisfy $\Con(e)\ge\tau$, Counter is shrunk by
\texttt{max\_C\_safe} to protect minimum Active evidence where possible.

\subsection{Toy numeric walk-through}
Suppose:
\begin{equation}
|\Pool|=9,\;\texttt{min\_A}=5,\;\texttt{max\_A}=20,\;\texttt{max\_C}=5.
\end{equation}
Then:
\begin{equation}
\texttt{max\_C\_safe}=\max(0,9-5)=4,
\end{equation}
so at most 4 items may enter Counter even if 5+ satisfy threshold.
Remaining 5 items fill Active exactly, preserving the starvation guard.
